\Article{236}
\textsc{Aufgabe.} \textit{Gegeben sind zwei Formen, deren Determinanten entweder gleich sind oder wenigstens im Verhältnis von Quadraten zueinander stehen. Gesucht ist eine aus ihnen zusammengesetzte Form.}

\textit{Lösung.} Die zusammenzusetzenden Formen seien $(a,b,c)\ldots f$, $(a',b',c')\ldots f'$; ihre Determinanten seien $d,d'$; die größten gemeinsamen Teiler der Zahlen $a,\allowbreak 2b,\allowbreak c;\allowbreak \ a',\allowbreak 2b',\allowbreak c'$ seien jeweils $m,m'$; der größte gemeinsame Teiler der Zahlen $dm'm',d'mm$, mit demselben Vorzeichen wie $d,d'$
\SourceLeaf{254}
versehen, sei $D$. Dann sind $\frac{dm'm'}{D},\frac{d'mm}{D}$ positive, zueinander teilerfremde Zahlen, deren Produkt ein Quadrat ist; also sind sie selbst Quadrate (Artikel 21). Damit sind $\sqrt{\frac{d}{D}},\sqrt{\frac{d'}{D}}$ rationale Größen, die wir $=n,n'$ setzen. Dabei nehmen wir für $n$ den positiven oder negativen Wert, je nachdem, ob die Form $f$ direkt oder invers in die Komposition eingehen soll; ebenso bestimmen wir das Vorzeichen von $n'$ danach, wie $f'$ in die Komposition eingehen soll. Also sind $mn',m'n$ zueinander teilerfremde ganze Zahlen; $n$ und $n'$ können aber auch Brüche sein. Nun bemerken wir, dass $an',cn',a'n,c'n,bn'+b'n,bn'-b'n$ ganze Zahlen sind. Für die ersten vier ist das unmittelbar klar (weil $an'=\frac{a}{m}mn'$ usw.); für die übrigen beiden beweist man es ebenso, wie im vorigen Artikel gezeigt wurde, dass $R$ und $S$ durch $e$ teilbar sind.

Nun wählen wir vier beliebige ganze Zahlen $\mathfrak{Q},\mathfrak{Q}',\mathfrak{Q}'',\mathfrak{Q}'''$, nur unter der Bedingung, dass die vier Größen auf der linken Seite der folgenden Gleichung (I) nicht alle zugleich $=0$ werden, und setzen
\[
\begin{aligned}
\mathfrak{Q}'an'+\mathfrak{Q}''a'n+\mathfrak{Q}'''(bn'+b'n)&=\mu q\qquad\ldotp\,\ldotp\,\ldotp\,\ldotp\,\ldotp\qquad\text{(I)}\\
-\mathfrak{Q}an'+\mathfrak{Q}'''c'n-\mathfrak{Q}''(bn'-b'n)&=\mu q'\\
\mathfrak{Q}'''cn'-\mathfrak{Q}a'n+\mathfrak{Q}'(bn'-b'n)&=\mu q''\\
-\mathfrak{Q}''cn'-\mathfrak{Q}'c'n-\mathfrak{Q}(bn'+b'n)&=\mu q'''
\end{aligned}
\]
so, dass $q,q',q'',q'''$ ganze Zahlen ohne gemeinsamen Teiler werden. Das erreicht man, indem man für $\mu$ den größten gemeinsamen Teiler der vier Zahlen auf den linken Seiten dieser Gleichungen nimmt. Dann lassen sich nach Artikel 40 vier ganze Zahlen $\mathfrak{P},\mathfrak{P}',\mathfrak{P}'',\mathfrak{P}'''$ finden, sodass
\[
\mathfrak{P}q+\mathfrak{P}'q'+\mathfrak{P}''q''+\mathfrak{P}'''q'''=1
\]
gilt. Danach bestimmen wir die Zahlen $p,p',p'',p'''$ durch die folgenden Gleichungen:
\[
\begin{aligned}
\mathfrak{P}'an'+\mathfrak{P}''a'n+\mathfrak{P}'''(bn'+b'n)&=p\qquad\ldotp\,\ldotp\,\ldotp\,\ldotp\,\ldotp\qquad\text{(II)}\\
-\mathfrak{P}an'+\mathfrak{P}'''c'n-\mathfrak{P}''(bn'-b'n)&=p'\\
\mathfrak{P}'''cn'-\mathfrak{P}a'n+\mathfrak{P}'(bn'-b'n)&=p''\\
-\mathfrak{P}''cn'-\mathfrak{P}'c'n-\mathfrak{P}(bn'+b'n)&=p'''
\end{aligned}
\]
Schließlich setzen wir
\[
\begin{gathered}q'q''-qq'''=Ann',\\pq'''+qp'''-p'q''-q'p''=2Bnn',\\p'p''-pp'''=Cnn'\end{gathered}
\]
\SourceLeaf{255}
Dann sind $A,B,C$ ganze Zahlen, und die Form $(A,B,C)\ldots F$ ist aus den Formen $f,f'$ zusammengesetzt.

\textit{Beweis.} I. Aus den Gleichungen I bestätigt man ohne Schwierigkeit die folgenden vier Gleichungen:
\[
\begin{aligned}
0&=q'cn'-q''c'n-q'''(bn'-b'n)\qquad\ldotp\,\ldotp\,\ldotp\,\ldotp\,\ldotp\,\ldotp\qquad\text{(III)}\\
0&=qcn'+q'''a'n-q''(bn'+b'n)\\
0&=q'''an'+qc'n-q'(bn'+b'n)\\
0&=q''an'-q'a'n-q(bn'-b'n)
\end{aligned}
\]

II. Nehmen wir nun an, die ganzen Zahlen $\mathfrak{A},\mathfrak{B},\mathfrak{C},\mathfrak{A}',\mathfrak{B}',\mathfrak{C}',\mathfrak{N},\mathfrak{N}'$ seien so bestimmt, dass
\[
\begin{aligned}
\mathfrak{A}a+2\mathfrak{B}b+\mathfrak{C}c&=m\\
\mathfrak{A}'a'+2\mathfrak{B}'b'+\mathfrak{C}'c'&=m'\\
\mathfrak{N}m'n+\mathfrak{N}'mn'&=1
\end{aligned}
\]
gilt. Dann ist
\[
\begin{aligned}\mathfrak{A}a\mathfrak{N}'n'&+2\mathfrak{B}b\mathfrak{N}'n'+\mathfrak{C}c\mathfrak{N}'n'\\&+\mathfrak{A}'a'\mathfrak{N}n+2\mathfrak{B}'b'\mathfrak{N}n+\mathfrak{C}'c'\mathfrak{N}n=1\end{aligned}
\]
Daraus und aus den Gleichungen (III) bestätigt man leicht: Setzt man
\[
\begin{aligned}
-q'\mathfrak{A}\mathfrak{N}'-q''\mathfrak{A}'\mathfrak{N}-q'''(\mathfrak{B}\mathfrak{N}'+\mathfrak{B}'\mathfrak{N})&=\mathfrak{q}\\
q\mathfrak{A}\mathfrak{N}'-q'''\mathfrak{C}'\mathfrak{N}+q''(\mathfrak{B}\mathfrak{N}'-\mathfrak{B}'\mathfrak{N})&=\mathfrak{q}'\\
-q'''\mathfrak{C}\mathfrak{N}'+q\mathfrak{A}'\mathfrak{N}-q'(\mathfrak{B}\mathfrak{N}'-\mathfrak{B}'\mathfrak{N})&=\mathfrak{q}''\\
q''\mathfrak{C}\mathfrak{N}'+q'\mathfrak{C}'\mathfrak{N}+q(\mathfrak{B}\mathfrak{N}'+\mathfrak{B}'\mathfrak{N})&=\mathfrak{q}'''
\end{aligned}
\]
so ist
\[
\begin{aligned}
\mathfrak{q}'an'+\mathfrak{q}''a'n+\mathfrak{q}'''(bn'+b'n)&=q\qquad\ldotp\,\ldotp\,\ldotp\,\ldotp\,\ldotp\,\ldotp\qquad\text{(IV)}\\
-\mathfrak{q}an'+\mathfrak{q}'''c'n-\mathfrak{q}''(bn'-b'n)&=q'\\
\mathfrak{q}'''cn'-\mathfrak{q}a'n+\mathfrak{q}'(bn'-b'n)&=q''\\
-\mathfrak{q}''cn'-\mathfrak{q}'c'n-\mathfrak{q}(bn'+b'n)&=q'''
\end{aligned}
\]

Wenn $\mu=1$ gilt, sind diese Gleichungen nicht nötig; an ihrer Stelle kann man die Gleichungen (I) beibehalten, denen sie völlig entsprechen. Entwickelt man nun aus den Gleichungen II, IV die Werte von $Ann',2Bnn',Cnn'$ (\textit{d.h.}, der Zahlen $q'q''-qq'''$ usw.) und streicht die Glieder, die einander aufheben, so findet man, dass die Glieder jeder dieser Größen
\SourceLeaf{256}
Produkte aus ganzen Zahlen mit $nn'$, mit $dn'n'$ oder mit $d'nn$ sind und außerdem alle Glieder von $2Bnn'$ den Faktor 2 enthalten. Daraus folgt (weil $dn'n'=d'nn$ und daher $\frac{dn'n'}{nn'}=\frac{d'nn}{nn'}=\sqrt{dd'}$ ganze Zahlen sind), dass $A,B,C$ ganze Zahlen sind. \textit{Q. E. P.}

III. Setzt man aus den Gleichungen (II) die Werte von $p,p',p'',p'''$ ein, so bestätigt man mithilfe der Gleichungen (III) und der Gleichung
\[
\mathfrak{P}q+\mathfrak{P}'q'+\mathfrak{P}''q''+\mathfrak{P}'''q'''=1
\]
leicht, dass
\[
\begin{aligned}
pq'-qp'&=an',\\ pq'''-qp'''-p'q''+q'p''&=2bn',\\p''q'''-q''p'''&=cn'\\
pq''-qp''&=a'n,\\ pq'''-qp'''+p'q''-q'p''&=2b'n,\\p'q'''-q'p'''&=c'n
\end{aligned}
\]
gilt. Diese Gleichungen sind mit den ersten sechs ($\Omega$) des vorigen Artikels identisch; die übrigen drei gelten bereits nach Voraussetzung. Also geht (\textit{am Ende desselben Artikels}) die Form $F$ in $ff'$ durch die Substitution $p,p',p'',p''';\ q,q',q''\,q'''$ über; ihre Determinante ist $=D$, also gleich dem größten gemeinsamen Teiler der Zahlen $dm'm',d'mm$, weshalb nach dem vierten Schluss des vorigen Artikels $F$ aus $f,f'$ zusammengesetzt ist. \textit{Q. E. S.} Schließlich sieht man leicht, dass $F$ aus $f,f'$ \textit{so} zusammengesetzt ist, wie es verlangt war, weil die Vorzeichen der Größen $n,n'$ schon von Anfang an richtig bestimmt wurden.
